% PDF 1.4 without object streams, for older viewers.
\pdfminorversion=4
\pdfobjcompresslevel=0
\documentclass[11pt]{article}
\usepackage[margin=2.3cm]{geometry}
\usepackage{amsmath,amssymb}
\usepackage{booktabs,array,longtable}
\usepackage[T1]{fontenc}
\usepackage{lmodern}
\usepackage{microtype}
\usepackage{tikz}
\usepackage[hidelinks]{hyperref}

\renewcommand{\P}{\mathcal{P}}
\newcommand{\F}{\mathbb{F}_3}
\newcommand{\A}{\mathcal{A}}
\DeclareMathOperator{\LT}{LT}
\DeclareMathOperator{\Ext}{Ext}
\setlength{\parindent}{0pt}
\setlength{\parskip}{0.4em}

\title{A minimal resolution over the Steenrod algebra at $p = 3$,\\ computed by Gröbner bases}
\author{Trace of the \texttt{groebner} crate (branch \texttt{groebner\_lin})}
\date{}

\begin{document}
\maketitle

\begin{abstract}
@ABSTRACT_LEAD@ the computation of the minimal free resolution
$\dots \to C_3 \to C_2 \to C_1 \to C_0 \to \F$ of the ground field over the mod $3$ Steenrod
algebra $\A$ by the \texttt{groebner} crate, a port of Weinan Lin's Gröbner basis method
extended to odd primes. The range is homological degree $s \le @S_MAX@$ and stem $t - s \le @N_MAX@$, so
internal degrees $t \le @T_MAX@$. Compared with $p = 2$, there are three new features: the exterior
generators $Q_k$, coefficients in $\F$, and polynomial generators of height $3$ in the
associated graded algebra. Every formula below is taken from the program's own data: the
critical pairs it formed, the elements it reduced, and the bases and differentials it
produced. Only the bookkeeping of the reduction steps was replayed outside the program, and the
replay was checked to reproduce the program's results exactly.
\end{abstract}

\section{Setting}

\paragraph{The algebra.}
We work in the Milnor basis $Q(E)\,\P(R) = Q_0^{\epsilon_0}Q_1^{\epsilon_1}\cdots
\P(r_1, r_2, \dots)$ of $\A$, writing $\P^r = \P(r)$. Here $Q_0 = \beta$ is the Bockstein,
$|Q_k| = 2\cdot 3^k - 1$, and $|\P(R)| = \sum_n 2 r_n (3^n - 1)$. Thus $|Q_0| = 1$, $|Q_1| = 5$,
$|\P^1| = 4$, and $|\P(0,1)| = 16$. Write $r_n = \sum_i d_{i,n} 3^i$ with ternary digits
$d_{i,n} \in \{0, 1, 2\}$. The digit records the power of $\xi_n^{3^i}$ in the dual basis,
of degree $2\cdot 3^i(3^n - 1)$. We record $Q(E)\,\P(R)$ by the monomial
\[ Q(E) \prod x_{in}^{d_{i,n}} \]
in the associated graded algebra of the May filtration
\[ E^0\A = E[Q_0, Q_1, \dots] \otimes \bigotimes_{i \ge 0,\, n \ge 1} \F[x_{in}]/(x_{in}^3). \]
For example $\P^2 \leftrightarrow x_{01}^2$, $\P^3 \leftrightarrow x_{11}$ and
$Q_0\P^1 \leftrightarrow Q_0x_{01}$. A monomial is given by a set of $Q_k$ and a digit
$0 \le d \le 2$ in each position. Divisibility is inclusion of the sets together with digitwise
$\le$; quotients, greatest common divisors and least common multiples are digitwise too. The
product in $E^0\A$ is non-zero exactly when no $Q_k$ repeats and no digit sum reaches $3$. Even
then it carries a coefficient: a binomial coefficient for each digit, and a sign from
reordering the $Q_k$. The positions in this range, in increasing degree, are @POSITIONS@.

\paragraph{The monomial order.}
A term of a free module is $c\, Q(E)\P(R)\, g$ with $c \in \F^\times$ and $g$ a generator.
Terms are compared by
\begin{enumerate}
\item the generator: the generator added later (of higher degree) comes first;
\item the May weight $\sum_{k \in E}(2k+1) + \sum d_{i,n}(2n - 1)$: smaller comes first;
\item the exterior part, read as a binary number with $Q_0$ most significant: smaller comes
  first;
\item the digits, read as a number with the position of lowest degree most significant:
  smaller comes first.
\end{enumerate}
The \emph{leading term} $\LT(x)$ is the first term in this order, and elements are always
printed with it first. The leading term of a product $Q(E)\P(R)\cdot x$ is the product of
$Q(E)\,x^R$ with $\LT(x)$ in $E^0\A$, \emph{including its coefficient}, whenever that product
is non-zero. Elements of a Gröbner basis are normalized to leading coefficient $1$. Since
$p = 3$, coefficients are printed as $\pm 1$.

\paragraph{Names.}
The generator of $C_0$ is $\iota$. The generators of $C_s$, $s \ge 1$, are named after the
classes of $\Ext_\A^{s,t}(\F, \F)$ they represent. The names of the decomposable classes are
found from the differentials: in a minimal resolution, the product $a_0 x$ or $h_i x$ of the
class $x$ dual to a generator $g$ of $C_s$ is non-zero on a generator $g'$ of $C_{s+1}$ exactly
when $Q_0 g$ or $\P^{3^i} g$ occurs in $d g'$. A generator is named $a_0 x$ or $h_i x$ when that
product is dual to it alone, up to sign. @NAME_NOTE@
These are only labels. The program numbers the generators of each $C_s$ in the order they are
found.

\section{The algorithm}

The algorithm is the same as at $p = 2$. We recall it briefly, with the changes for odd primes.

\paragraph{Levels.}
Level $s$ is a Gröbner basis $b^s_0, b^s_1, \dots$ of the graph of the differential,
\[ \Gamma_s = \{ (d y, y) : y \in C_{s+1} \} \subseteq C_s \oplus C_{s+1}, \]
for the order on the first component. Each $b^s_k = (x_1 \mid x_2)$ satisfies $d x_2 = x_1$, its
leading term is that of $x_1$, and its leading coefficient is $1$. The module being resolved
enters as ``level $-1$'': $\F = \A/\A\{Q_0, \P^1, \P^3, \P^9, \dots\}$, so the relations are
$Q_0\iota = 0$ and $\P^{3^i}\iota = 0$.

\paragraph{Critical pairs and singles.}
When an element with leading term $L$ joins level $s$, the program forms
\begin{itemize}
\item a \emph{critical pair} with each earlier element $b$ whose leading term $L'$ lies on the
  same generator. Let $M = \operatorname{lcm}(L, L')$ in $E^0\A$, with cofactors $u = M/L'$ and
  $u' = M/L$. The products $u\cdot b$ and $u'\cdot b_{\text{new}}$ have leading terms
  $\lambda M$ and $\mu M$. The S-element is $\lambda\, u'\cdot b_{\text{new}} - \mu\, u\cdot b$,
  in which the leading terms cancel. The pairs are pruned by the Gebauer--Möller criteria;
\item a \emph{critical single} $m \cdot b_{\text{new}}$ for each generator $m$ of the
  annihilator of $L$ in $E^0\A$. These are $m = x_{in}^{3-d}$ for each digit $d$ of $L$ in
  position $x_{in}$, and $m = Q_k$ for each $Q_k$ in $L$. Then $m \cdot L = 0$ in $E^0\A$, so
  the product must be reduced.
\end{itemize}

\paragraph{One degree.}
In degree $t$, the program does the following:
\begin{enumerate}
\item Every level hands out its pairs and singles of degree $t$, and the relations of degree
  $t$ are taken up.
\item Each of them is reduced. The first component is reduced modulo level $s$: a term
  $c\,m$ divisible by $\LT(b)$, with quotient $u$, is cancelled by adding $-c/\lambda$ times
  $u \cdot b$, where $\lambda$ is the leading coefficient of $u \cdot b$. The second component
  is carried along. Then the second component is reduced modulo the first components of level
  $s+1$, i.e.\ modulo the boundaries known so far. If the first component has become zero, the
  second component $z \in C_{s+1}$ is a cycle.
\item Going down from the top level, the results are compared with the elements found earlier
  in the same degree, normalized, and sorted. A non-zero first component gives a \textbf{new
  basis element}. If both components are zero, the reduction was \textbf{useless}. Otherwise
  $z$ is a cycle that is not a boundary, and it becomes $d$ of a \textbf{new generator} of
  $C_{s+2}$, or the outcome is \textbf{beyond} if $s + 2 > @S_MAX@$.
\end{enumerate}

\paragraph{Truncation.}
We want stems $\le @N_MAX@$. A cycle found at level $s$ in degree $t$ is a generator of $C_{s+2}$ in
stem $t-s-2$, so level $s$ is only needed for $t \le s + @N_MAX2@$. Relations are only needed for
$t \le @N_MAX1@$.

\paragraph{Running it.}\sloppy
The computation was run by a small driver around the public API of the crate, in
\texttt{docs/trace}. It runs \texttt{Resolution::sphere} at $p = 3$ with
\texttt{SteenrodKind::Classical} and \texttt{Bounds \{ s\_max: @S_MAX@, n\_max: @N_MAX@ \}}, and
calls \texttt{compute\_through(@T_MAX@)}. It also replays the divisibility criterion on the final
leading terms, and checks that it prunes as many pairs in each degree as the program counted.

@STEPS@
\section{The result}

@RESULT@
